\documentclass[11pt,reqno]{amsart}

\usepackage[T1]{fontenc}
\usepackage[utf8]{inputenc}
\usepackage{newtxtext,newtxmath}
\usepackage{microtype}
\usepackage{amsmath,mathtools}
\usepackage{enumitem}
\usepackage[hidelinks]{hyperref}
\usepackage[nameinlink,noabbrev]{cleveref}

\newcommand{\R}{\mathbf R}
\newcommand{\C}{\mathbf C}
\newcommand{\PP}{\mathbf P}
\newcommand{\Gm}{\mathbf G_m}
\newcommand{\NOne}[1]{N^1(#1)_{\R}}
\newcommand{\OO}{\mathcal O}
\DeclareMathOperator{\Aut}{Aut}
\DeclareMathOperator{\Bl}{Bl}
\DeclareMathOperator{\Hess}{Hess}
\DeclareMathOperator{\ord}{ord}
\DeclareMathOperator{\lcm}{lcm}

\newtheorem{theorem}{Theorem}[section]
\newtheorem{proposition}[theorem]{Proposition}
\newtheorem{lemma}[theorem]{Lemma}
\newtheorem{corollary}[theorem]{Corollary}
\theoremstyle{remark}
\newtheorem{remark}[theorem]{Remark}

\title[Hessian discriminants and blow-ups]{Hessian discriminants and automorphisms of blow-ups}
\author{Open AI}
\date{}

\subjclass[2020]{14J50, 14E05, 14C17}
\keywords{blow-up, automorphism group, intersection form, Hessian discriminant, exceptional divisor}

\begin{document}

\begin{abstract}
Let $\pi\colon Y=\Bl_S X\to X$ be the blow-up of a smooth complex projective $n$-fold along a nonempty connected smooth centre of codimension $c\ge2$.  Write $W=\pi^*N^1(X)_{\R}\subset N^1(Y)_{\R}$ and let $\Delta_Y$ be the determinant of the Hessian bilinear form of the top intersection polynomial.  We prove that $W$ occurs in the Hessian discriminant with exact multiplicity $c-2$.  For $c\ge3$, it is therefore a linear component with finite orbit under $\Aut(Y)$.  Its stabilizer is exactly the subgroup obtained by lifting $\Aut(X,S)$, and
\[
 [\Aut(Y):\Aut(X,S)]
 \le
 \left\lfloor\frac{(n-2)\rho(Y)}{c-2}\right\rfloor.
\]
For a finite sequence of blow-ups along smooth centres of codimension at least three, a uniform power of every automorphism descends through the whole sequence.  Consequently, if the initial smooth complex projective variety has finite automorphism group, then so does the final blow-up.  Thus the examples asked for in Lesieutre--Litt's Question~1.8 do not occur when the initial variety is smooth over $\C$.
\end{abstract}

\maketitle

\section{Introduction}

Let $X$ be a smooth complex projective variety and let $Y$ be obtained from $X$ by a sequence of ordinary blow-ups along smooth centres.  A basic question is whether the birational modification can create new automorphisms.  Bayraktar--Cantat obtained strong restrictions when the centres have sufficiently small dimension, and related questions were studied by Truong; see \cite{BayraktarCantat2013,Truong2013}.  Lesieutre--Litt proved a uniform descent theorem under the hypotheses
\[
\begin{aligned}
&2r+3\le n,\\
&\text{or } r+3\le n\text{ together with polyhedrality of the exceptional nef cone}.
\end{aligned}
\]
for a centre of dimension $r$ in an $n$-fold \cite[Theorem~1.6]{LesieutreLitt2019}.  They left open the intermediate range and asked whether blow-ups along smooth centres of codimension at least three can turn a finite automorphism group into an infinite one \cite[Question~1.8]{LesieutreLitt2019}.

The argument below uses only the top intersection polynomial on $N^1$.  For a smooth projective $n$-fold $Z$, set
\[
 P_Z(h)=\int_Z h^n,
 \qquad
 q_h(u,v)=\int_Z h^{n-2}uv,
 \qquad
 \Delta_Z(h)=\det(q_h).
\]
Up to a nonzero scalar, $\Delta_Z$ is the determinant of the ordinary Hessian of $P_Z$.  If $\pi\colon Y=\Bl_SX\to X$ has centre of codimension $c$, then the pullback space
\[
 W=\pi^*N^1(X)_{\R}
\]
is a hyperplane in $N^1(Y)_{\R}$.  The key point is the exact factorization order
\begin{equation}\label{eq:intro-multiplicity}
 \ord_W(\Delta_Y)=c-2.
\end{equation}
Thus codimension three is exactly the first codimension for which the blow-down is recorded by a linear component of the Hessian discriminant.  Automorphisms preserve $\Delta_Y$ up to a nonzero scalar, hence permute these linear components with their multiplicities.  The orbit of $W$ is finite, and its stabilizer turns out to be exactly the subgroup lifted from automorphisms of the pair $(X,S)$.

Intersection forms have long been used to constrain birational geometry.  In dimension three, Bisi--Cascini--Tasin used the cubic intersection form to obstruct smooth blow-up presentations \cite{BisiCasciniTasin2016}.  More recently, they studied the hypermatrix of the $(n-1)$st derivatives of the full intersection form; in particular, they derived a blow-up expansion and used low-rank exceptional classes to constrain divisorial contractions \cite[Sections~4--5]{BisiCasciniTasin2025}.  The invariant used here is different: it is the determinant of the ordinary Hessian on $N^1$, and the relevant locus is the pullback hyperplane $W$, rather than the exceptional class itself.

The exact multiplicity in \eqref{eq:intro-multiplicity} yields both an explicit index bound and a uniform iterate theorem.  The latter removes the dimension and nef-cone hypotheses from Lesieutre--Litt's descent theorem when every centre has codimension at least three.

\section{The Hessian discriminant}

Let $Z$ be a smooth complex projective $n$-fold, $n\ge3$, and put
\[
 V_Z=N^1(Z)_{\R},\qquad \rho(Z)=\dim_{\R}V_Z.
\]
Choose a basis $D_1,\ldots,D_\rho$ of $V_Z$.  Define
\begin{equation}\label{eq:Delta-def}
 q_h(u,v)=\int_Z h^{n-2}uv,
 \qquad
 M_Z(h)=\bigl(q_h(D_i,D_j)\bigr)_{i,j},
 \qquad
 \Delta_Z(h)=\det M_Z(h).
\end{equation}
If $h=\sum x_iD_i$, then
\[
 \det\Hess(P_Z)
 =\bigl(n(n-1)\bigr)^\rho\Delta_Z.
\]
A change of basis multiplies $\Delta_Z$ by a nonzero square.  Hence the zero hypersurface of $\Delta_Z$, and the multiplicity of each real linear factor, are intrinsic.

\begin{lemma}\label{lem:discriminant}
The polynomial $\Delta_Z$ is nonzero and homogeneous of degree $(n-2)\rho(Z)$.  Its zero hypersurface is preserved by $\Aut(Z)$, including the multiplicity of every irreducible factor.
\end{lemma}

\begin{proof}
For an ample class $a$, the Hodge index theorem gives signature $(1,\rho(Z)-1)$ for
\[
 (u,v)\longmapsto\int_Z a^{n-2}uv
\]
on $N^1(Z)_{\R}$; see, for example, \cite[Theorem~7.3.4]{deCataldo2005}.  Hence $\Delta_Z(a)\ne0$.  Each entry of $M_Z(h)$ is homogeneous of degree $n-2$, so $\deg\Delta_Z=(n-2)\rho(Z)$.

If $g\in\Aut(Z)$ and $T$ is the matrix of $g^*$ in the chosen basis, invariance of intersection numbers gives
\[
 T^{\mathsf t}M_Z(g^*h)T=M_Z(h).
\]
Thus
\[
 \Delta_Z(g^*h)=(\det T)^{-2}\Delta_Z(h).
\]
Pullback therefore preserves the zero hypersurface of $\Delta_Z$ and all factor multiplicities.
\end{proof}

\section{The multiplicity of the blow-down hyperplane}

Let
\[
 \pi\colon Y=\Bl_SX\longrightarrow X
\]
be the blow-up of a smooth complex projective $n$-fold along a nonempty connected smooth centre $S$ of codimension $c\ge2$.  Let $E$ be the exceptional divisor.  Then
\begin{equation}\label{eq:NS-decomp}
 N^1(Y)_{\R}=\pi^*N^1(X)_{\R}\oplus\R[E].
\end{equation}
Write
\[
 W=\pi^*N^1(X)_{\R}.
\]

We first record the only intersection calculation needed below.

\begin{lemma}\label{lem:push-E}
For $1\le j<c$ one has $\pi_*E^j=0$, while
\[
 \pi_*E^c=(-1)^{c-1}[S].
\]
\end{lemma}

\begin{proof}
Let $p\colon E=\PP(N_{S/X})\to S$ and put $\xi=c_1(\OO_E(1))$.  Since $\OO_E(E)=\OO_E(-1)$,
\[
 E^j=i_*((-\xi)^{j-1}).
\]
The projective bundle formula gives $p_*\xi^k=0$ for $k<c-1$ and $p_*\xi^{c-1}=1$.  Pushing forward through $S\hookrightarrow X$ proves the claim; compare \cite[Chapter~3]{Fulton1998}.
\end{proof}

\begin{theorem}[Exact Hessian multiplicity]\label{thm:multiplicity}
With the notation above,
\[
 \boxed{\ord_W(\Delta_Y)=c-2.}
\]
Equivalently, if $\ell_W$ is a linear equation for $W$, then
\[
 \ell_W^{\,c-2}\mid\Delta_Y,
 \qquad
 \ell_W^{\,c-1}\nmid\Delta_Y.
\]
In particular, $W$ is a component of the Hessian discriminant if and only if $c\ge3$.
\end{theorem}

\begin{proof}
Choose a basis $D_1,\ldots,D_r$ of $N^1(X)_{\R}$ and use
\[
 \pi^*D_1,\ldots,\pi^*D_r,E
\]
as a basis of $N^1(Y)_{\R}$.  Write
\[
 h=\pi^*a+tE.
\]
Relative to this basis, the matrix of $q_h$ has block form
\[
 M_Y(h)=
 \begin{pmatrix}
   A(t)&b(t)\\
   b(t)^{\mathsf t}&d(t)
 \end{pmatrix}.
\]
By \cref{lem:push-E} and the projection formula,
\begin{align}
 A(t)&=A(0)+O(t^c),\label{eq:Aorder}\\
 b(t)&=O(t^{c-1}),\label{eq:border}\\
 d(t)&=(-1)^{c-1}\binom{n-2}{c-2}
       \left(\int_S(a|_S)^{n-c}\right)t^{c-2}
       +O(t^{c-1}).\label{eq:dorder}
\end{align}
The block estimates hold in the polynomial ring in the coordinates of $a$ and $t$.  Expanding the determinant by permutations therefore shows directly that
\[
 t^{c-2}\mid \det M_Y(\pi^*a+tE)
\]
as a polynomial, not merely after specializing $a$.

Now take $a$ ample.  The Hodge index theorem gives $\det A(0)\ne0$, and ampleness on $S$ gives
\[
 \int_S(a|_S)^{n-c}>0.
\]
In the formal power series ring at $t=0$,
\[
 \det M_Y(h)
 =\det A(t)\bigl(d(t)-b(t)^{\mathsf t}A(t)^{-1}b(t)\bigr).
\]
The second term in the parentheses has order at least $2c-2$, whereas \eqref{eq:dorder} has a nonzero term of order $c-2$.  Therefore the coefficient of $t^{c-2}$ in $\det M_Y(h)$ is nonzero for this ample $a$.  The global divisibility by $t^{c-2}$ is consequently exact.
\end{proof}

\begin{remark}[The codimension-two boundary]\label{rem:c2}
For $c=2$, \cref{thm:multiplicity} gives multiplicity zero.  More concretely, for ample $a$ the block $A(0)$ is nondegenerate, the mixed block vanishes at $t=0$, and
\[
 d(0)=-\int_S(a|_S)^{n-2}<0.
\]
Hence $\Delta_Y(\pi^*a)\ne0$.  This is exactly where the Hessian-hyperplane mechanism stops.  Lesieutre--Litt exhibit codimension-two blow-ups for which uniform descent fails \cite[Example~2.5]{LesieutreLitt2019}.
\end{remark}

\section{Automorphisms of a single blow-up}

Write
\[
 \Aut(X,S)=\{\varphi\in\Aut(X):\varphi(S)=S\}.
\]
Every element of $\Aut(X,S)$ lifts uniquely to $Y$.

\begin{theorem}[Virtual recovery of the blow-down]\label{thm:single}
Assume $c\ge3$ and set
\[
 M(Y,S)=
 \left\lfloor
 \frac{(n-2)\rho(Y)}{c-2}
 \right\rfloor.
\]
Then the natural lift identifies $\Aut(X,S)$ with the stabilizer of $W$ in $\Aut(Y)$.  In particular,
\begin{equation}\label{eq:index-bound}
 [\Aut(Y):\Aut(X,S)]\le M(Y,S).
\end{equation}
Moreover, if
\[
 L(Y,S)=\lcm(1,2,\ldots,M(Y,S)),
\]
then for every $g\in\Aut(Y)$ the iterate $g^{L(Y,S)}$ descends to an automorphism of $X$ preserving $S$.
\end{theorem}

\begin{proof}
By \cref{thm:multiplicity}, the linear form defining $W$ divides $\Delta_Y$ with multiplicity $c-2$.  By \cref{lem:discriminant}, $\Aut(Y)$ permutes the linear factors having this multiplicity.  There are at most
\[
 \left\lfloor\frac{\deg\Delta_Y}{c-2}\right\rfloor
 =M(Y,S).
\]
Thus the stabilizer
\[
 G_W=\{g\in\Aut(Y):g^*W=W\}
\]
has index at most $M(Y,S)$.

It remains to identify $G_W$.  Let $g\in G_W$ and choose an ample divisor $A$ on $X$.  If $C\subset Y$ is an integral curve contracted by $\pi$, then
\[
 A\cdot\pi_*g_*[C]
 =(g^*\pi^*A)\cdot C
 =0,
\]
because $g^*\pi^*A\in W$.  If $\pi(g(C))$ were a curve, ampleness would make the left side positive.  Hence $g$ sends every contracted curve to a contracted curve.

Each nontrivial fibre of $\pi$ is $\PP^{c-1}$, whose closed points are joined by lines.  Therefore $\pi g$ is constant on every fibre.  Let $\Gamma$ be the reduced image of
\[
 (\pi,\pi g)\colon Y\longrightarrow X\times X.
\]
This map is proper, so its image is closed.  The first projection $q\colon\Gamma\to X$ is proper.  For every $x\in X$, all points of $\pi^{-1}(x)$ have the same image in $\Gamma$; since $Y\to\Gamma$ is surjective, the underlying topological space of $q^{-1}(x)$ consists of one point.  Thus $q$ has finite fibres and is finite \cite[Tag~02LS]{Stacks}.  Over $X\setminus S$ it is the graph of the morphism $\pi g\pi^{-1}$, so $q$ is birational.  The variety $\Gamma$ is integral, and $X$ is normal; hence $q$ is an isomorphism \cite[Tag~0AB1]{Stacks}.  The second projection therefore gives a morphism $\bar g\colon X\to X$ satisfying
\[
 \pi g=\bar g\pi.
\]
Applying the same argument to $g^{-1}$ gives a morphism $\overline{g^{-1}}\colon X\to X$.  Since $\pi$ is surjective,
\[
 \bar g\,\overline{g^{-1}}=\overline{g^{-1}}\,\bar g=\mathrm{id}_X,
\]
so $\bar g\in\Aut(X)$.

For $y\in E$, choose a line in the fibre of $E\to S$ through $y$.  Its image under $g$ is contracted by $\pi$.  Every curve contracted by $\pi$ is contained in $E$, because $\pi$ is an isomorphism over $X\setminus S$.  Hence $g(E)\subseteq E$; since both are irreducible divisors, $g(E)=E$.  Therefore
\[
 \bar g(S)=\pi g(E)=S.
\]
If $\bar g=\mathrm{id}_X$, then $g$ is the identity on $Y\setminus E$, and hence on $Y$.  Therefore descent gives an injection $G_W\hookrightarrow\Aut(X,S)$.  Conversely, every automorphism of $(X,S)$ lifts uniquely to the blow-up by its universal property, and the lift preserves $W$.  Hence
\[
 G_W\cong\Aut(X,S),
\]
which proves \eqref{eq:index-bound}.

Finally, the orbit of $W$ under a fixed $g$ has length at most $M(Y,S)$.  Therefore $g^m$ stabilizes $W$ for some $1\le m\le M(Y,S)$, and $g^{L(Y,S)}$ does so uniformly for every $g$.
\end{proof}

\section{Sequences of blow-ups}

A blow-up along a finite disjoint union of smooth connected centres is the same as the successive blow-up of the components.  We therefore state the result for connected centres.

\begin{theorem}[Uniform descent]\label{thm:sequence}
Let
\[
 Y=X_m\xrightarrow{\pi_m}X_{m-1}\longrightarrow\cdots
 \xrightarrow{\pi_1}X_0=X
\]
be a finite sequence of smooth complex projective $n$-folds, where $\pi_i$ is the blow-up along a nonempty connected smooth centre $S_i\subset X_{i-1}$ of codimension $c_i\ge3$.  Put
\[
 M_i=
 \left\lfloor
 \frac{(n-2)\rho(X_i)}{c_i-2}
 \right\rfloor,
 \qquad
 L_i=\lcm(1,\ldots,M_i),
 \qquad
 N=\prod_{i=1}^m L_i.
\]
Then:
\begin{enumerate}[label=\textup{(\roman*)}]
\item there is a finite-index subgroup $G\le\Aut(Y)$ that descends faithfully to $\Aut(X)$;
\item
\[
 [\Aut(Y):G]\le\prod_{i=1}^m M_i;
\]
\item for every $g\in\Aut(Y)$, the iterate $g^N$ descends to an automorphism of $X$.
\end{enumerate}
Consequently, if $\Aut(X)$ is finite, then $\Aut(Y)$ is finite.
\end{theorem}

\begin{proof}
For each $i$, \cref{thm:single} gives a subgroup $G_i\le\Aut(X_i)$ of index at most $M_i$ together with an injective descent homomorphism
\[
 \delta_i\colon G_i\longrightarrow\Aut(X_{i-1}).
\]
Define recursively $H_0=\Aut(X_0)$ and
\[
 H_i=\delta_i^{-1}(H_{i-1})\subset G_i.
\]
Then
\[
 [G_i:H_i]
 =[\delta_i(G_i):\delta_i(G_i)\cap H_{i-1}]
 \le[\Aut(X_{i-1}):H_{i-1}].
\]
Induction gives
\[
 [\Aut(X_i):H_i]\le\prod_{j=1}^iM_j.
\]
Taking $G=H_m$ proves (i) and (ii), and faithfulness follows from the single-step faithfulness at every stage.

For (iii), start with $g_m=g$.  By \cref{thm:single}, $g_m^{L_m}$ descends to some $g_{m-1}\in\Aut(X_{m-1})$.  Then $g_{m-1}^{L_{m-1}}$ descends to $X_{m-2}$, and so on.  After $m$ steps, $g^N$ descends to $X$.  If $\Aut(X)$ is finite, then $G$ is finite by injectivity; since $G$ has finite index, so is $\Aut(Y)$.
\end{proof}

\begin{corollary}\label{cor:LL}
Let $X$ be smooth and complex projective.  For any finite sequence of blow-ups of $X$ along smooth centres of codimension at least three, Lesieutre--Litt's uniform-iterate conclusion holds with no additional dimension inequality and no polyhedral-nef-cone hypothesis.  In particular, no such sequence produces an infinite automorphism group from a finite one; this rules out the examples sought in \cite[Question~1.8]{LesieutreLitt2019} when the initial variety is smooth over $\C$.
\end{corollary}

\begin{remark}\label{rem:scope}
The proof is numerical except for the nondegeneracy of the Hodge-index form used in \cref{lem:discriminant}.  No dynamical Mordell--Lang input is required.  We have kept the statement over $\C$ in order to isolate the argument from additional characteristic-dependent hypotheses.
\end{remark}

\section*{AI assistance disclosure}
GPT-6 Astra and GPT-5.6 Sol were used during parts of the investigation, drafting, and checking of this note.

\begin{thebibliography}{99}

\bibitem{BayraktarCantat2013}
T.~Bayraktar and S.~Cantat,
\emph{Constraints on automorphism groups of higher dimensional manifolds},
J. Math. Anal. Appl. \textbf{405} (2013), no.~1, 209--213.

\bibitem{BisiCasciniTasin2016}
C.~Bisi, P.~Cascini, and L.~Tasin,
\emph{A remark on the Ueno--Campana's threefold},
Michigan Math. J. \textbf{65} (2016), no.~3, 567--572.

\bibitem{BisiCasciniTasin2025}
C.~Bisi, P.~Cascini, and L.~Tasin,
\emph{Forms and complex manifolds},
EMS Surv. Math. Sci. (2025), published online first, DOI: 10.4171/EMSS/109.

\bibitem{deCataldo2005}
M.~A.~A. de~Cataldo,
\emph{Lectures on the Hodge theory of projective manifolds},
arXiv:math/0504561 (2005).

\bibitem{Fulton1998}
W.~Fulton,
\emph{Intersection Theory}, 2nd ed.,
Springer, 1998.

\bibitem{LesieutreLitt2019}
J.~Lesieutre and D.~Litt,
\emph{Dynamical Mordell--Lang and automorphisms of blow-ups},
Algebraic Geometry \textbf{6} (2019), no.~1, 1--25.

\bibitem{Stacks}
The Stacks Project Authors,
\emph{The Stacks Project},
\url{https://stacks.math.columbia.edu}, Tags 02LS and 0AB1.

\bibitem{Truong2013}
T.~T. Truong,
\emph{On automorphisms of blowups of projective manifolds},
arXiv:1301.4957 (2013).

\end{thebibliography}

\end{document}
